\documentclass[12pt,a4paper]{article}

% =========================================================================
% FORMATAÇÃO ABNT (NBR 14724 / NBR 6024 / NBR 6027)
% =========================================================================

% --- Margens ABNT (3 cm superior e esquerda; 2 cm inferior e direita) ---
\usepackage[
  a4paper,
  left=3cm,
  top=3cm,
  right=2cm,
  bottom=2cm,
  headheight=15pt,
  headsep=1cm
]{geometry}

% --- Idioma, codificação e tipografia (Times New Roman - ABNT NBR 14724) ---
\usepackage[T1]{fontenc}
\usepackage[utf8]{inputenc}
\usepackage[brazil]{babel}
\usepackage{amsmath}
\usepackage{newtxtext,newtxmath}
\usepackage{microtype}
\usepackage[htt]{hyphenat}

% --- Espaçamento e parágrafo ABNT ---
\usepackage{setspace}
\onehalfspacing
\usepackage{indentfirst}
\setlength{\parindent}{1.25cm}
\setlength{\parskip}{0pt}
\setlength{\emergencystretch}{2em}

% --- Tabelas e alinhamentos ---
\usepackage{array}
\usepackage{tabularx}
\usepackage{booktabs}
\usepackage{etoolbox}

% --- Figuras, quadros e imagens ---
\usepackage{graphicx}
\graphicspath{{../assets/}{assets/}}
\usepackage{float}

% --- Cabeçalhos e paginação ABNT ---
\usepackage{fancyhdr}
\pagestyle{empty}
\fancyhf{}
\fancyhead[R]{\small\thepage}
\renewcommand{\headrulewidth}{0pt}

% --- Formatação do Sumário (NBR 6027) ---
\usepackage{tocloft}
\renewcommand{\cftsecleader}{\cftdotfill{\cftdotsep}}
\renewcommand{\cftsecfont}{\bfseries}
\renewcommand{\cftsecpagefont}{\bfseries}
\setlength{\cftbeforesecskip}{0.4em}

% --- Formatação dos títulos de seções (NBR 6024) ---
\usepackage{titlesec}
\titleformat{\section}{\normalfont\normalsize\bfseries\MakeUppercase}{\thesection}{1em}{}
\titleformat{\subsection}{\normalfont\normalsize\bfseries}{\thesubsection}{1em}{}
\titleformat{\subsubsection}{\normalfont\normalsize\bfseries}{\thesubsubsection}{1em}{}
\titleformat{\paragraph}{\normalfont\normalsize\itshape}{\theparagraph}{1em}{}

\titlespacing*{\section}{0pt}{2ex plus 0.5ex minus 0.2ex}{1.5ex plus 0.5ex minus 0.2ex}
\titlespacing*{\subsection}{0pt}{1.8ex plus 0.5ex minus 0.2ex}{1.2ex plus 0.5ex minus 0.2ex}
\titlespacing*{\subsubsection}{0pt}{1.5ex plus 0.5ex minus 0.2ex}{1.0ex plus 0.5ex minus 0.2ex}

% --- Hiperlinks ---
\usepackage{xcolor}
\definecolor{NavyBlue}{RGB}{0,40,120}
\usepackage[colorlinks=true,linkcolor=NavyBlue,citecolor=NavyBlue,urlcolor=NavyBlue]{hyperref}

\begin{document}

% -------------------------------------------------------------------------
% CAPA (NBR 14724)
% -------------------------------------------------------------------------
\begin{titlepage}
  \centering
  \MakeUppercase{\textbf{Universidade Federal de Goiás - UFG}}\\
  \MakeUppercase{\textbf{Instituto de Informática - INF}}\\
  \MakeUppercase{\textbf{Bacharelado em Engenharia de Software}}\\
  \vspace{3.2cm}
  
  \textbf{EQUIPE: CINESOFT SOLUTIONS}\\
  \vspace{0.6cm}
  \textbf{Carlos Eduardo Freire da Silva}\\
  \textbf{Gabriel Henrique Lourenço Dias}\\
  \textbf{Kayque Vicente Silva}\\
  \textbf{Yury Nogueira de Oliveira Silva}\\
  \vspace{3.2cm}
  
  \MakeUppercase{\textbf{Manual de Processos de Software}}\\
  \MakeUppercase{\textbf{Sistema CineDoom}}\\
  \vspace{0.8cm}
  \large Alinhado à Norma ABNT NBR ISO/IEC/IEEE 12207:2021\\
  \vfill
  
  \textbf{Goiânia - GO}\\
  \textbf{2026}
\end{titlepage}

% -------------------------------------------------------------------------
% FOLHA DE ROSTO
% -------------------------------------------------------------------------
\begin{titlepage}
  \centering
  \textbf{EQUIPE: CINESOFT SOLUTIONS}\\
  \vspace{0.4cm}
  \textbf{Carlos Eduardo Freire da Silva, Gabriel Henrique Lourenço Dias,}\\
  \textbf{Kayque Vicente Silva, Yury Nogueira de Oliveira Silva}\\
  \vspace{3.8cm}
  
  \MakeUppercase{\textbf{Manual de Processos de Software}}\\
  \MakeUppercase{\textbf{Sistema CineDoom}}\\
  \vspace{2.2cm}
  
  \hfill
  \begin{minipage}{0.52\textwidth}
    \begin{singlespace}
      \small
      Manual de Processos de Software elaborado para a disciplina de Processos de Engenharia de Software do curso de Bacharelado em Engenharia de Software do Instituto de Informática da Universidade Federal de Goiás (UFG), sob orientação da Profa. Dra. Juliana Pereira de Souza Zinader, em conformidade com o padrão ABNT NBR ISO/IEC/IEEE 12207:2021.
    \end{singlespace}
  \end{minipage}
  \vfill
  
  \textbf{Goiânia - GO}\\
  \textbf{2026}
\end{titlepage}

% -------------------------------------------------------------------------
% HISTÓRICO DE VERSÕES
% -------------------------------------------------------------------------
\newpage
\thispagestyle{empty}
\begin{center}
  \MakeUppercase{\textbf{Histórico de Versões do Documento}}
\end{center}
\vspace{0.8cm}

\noindent Documento controlado elaborado pela equipe \textbf{CineSoft Solutions} para o sistema \textbf{CineDoom}. Qualquer alteração neste manual é registrada formalmente nesta tabela e mantida em rastreabilidade com o catálogo formal de decisões de arquitetura (ADRs) e especificações do projeto.

\vspace{0.6cm}

\noindent
\begin{tabularx}{\textwidth}{@{}>{\centering\arraybackslash}p{2.5cm} >{\centering\arraybackslash}p{2cm} X@{}}
  \toprule
  \textbf{Data} & \textbf{Versão} & \textbf{Alteração Realizada} \\
  \midrule
  10/09/2026 & 1.0 & Criação do documento e elaboração da capa. \\
  \addlinespace[0.3em]
  16/09/2026 & 1.1 & Elaboração dos itens do Capítulo 1 (Apresentação do Software). \\
  \addlinespace[0.3em]
  24/09/2026 & 1.2 & Alinhamento dos Capítulos 1, 2 e 3 às especificações de engenharia e à norma ABNT NBR ISO/IEC/IEEE 12207:2021. \\
  \addlinespace[0.3em]
  24/09/2026 & 1.3 & Revisão técnica: correção de contradições, terminologia e índice de aderência. \\
  \addlinespace[0.3em]
  08/10/2026 & 1.4 & Incremento e estruturação dos processos do Capítulo 4 (Gerenciamento Técnico). \\
  \addlinespace[0.3em]
  17/10/2026 & 1.5 & Atualização da capa e padronização do histórico de versões conforme diretrizes docentes. \\
  \addlinespace[0.3em]
  22/10/2026 & 1.6 & Padronização formal da estrutura em Capítulos 1 a 5 conforme diretriz docente de estrutura mínima. \\
  \addlinespace[0.3em]
  23/10/2026 & 1.7 & Incorporação do Portal do Cliente (3º módulo), autenticação LGPD, reserva ACID, e-mail transacional e cancelamento web (ADR-012). \\
  \bottomrule
\end{tabularx}
\vfill
\newpage

% -------------------------------------------------------------------------
% SUMÁRIO
% -------------------------------------------------------------------------
\pagestyle{empty}
\tableofcontents
\newpage

% -------------------------------------------------------------------------
% CORPO DO DOCUMENTO
% -------------------------------------------------------------------------
\pagestyle{fancy}
\setcounter{page}{5}

\section{Capítulo 1 --- Apresentação do Software}
O sistema \textbf{CineDoom} é uma solução integrada de gestão operacional de salas de cinema e emissão de ingressos, concebida no modelo SaaS para pequenos e médios estabelecimentos. O presente manual estabelece os processos de ciclo de vida adotados conforme a norma ABNT NBR ISO/IEC/IEEE 12207:2021.

\subsection{Nome e Descrição do Software}
\subsection{Problema e Solução Proposta}
\subsection{Público-Alvo e Contexto de Uso}
\subsection{Usuários e Ambiente Operacional}
\subsection{Escopo e Funcionalidades Centrais}
\subsection{Restrições e Premissas}

\section{Capítulo 2 --- Processos de Acordo}
\subsection{Processo de Aquisição}
\subsection{Processo de Fornecimento}
\subsection{Fornecedores e Parceiros Relevantes}

\section{Capítulo 3 --- Processos Organizacionais Habilitadores}
\subsection{Modelo de Ciclo de Vida Adotado}
\subsection{Inventário de Infraestrutura e Ferramentas}
\subsection{Justificativa Estratégica e Viabilidade}
\subsection{Papéis, Competências e Matriz RACI}
\subsection{Plano de Qualidade e Métricas}
\subsection{Estratégia de Gestão do Conhecimento}

\section{Capítulo 4 --- Processos de Gerenciamento Técnico}
\subsection{Plano de Projeto, Escopo e Cronograma}
\subsection{Plano de Gestão de Riscos}
\subsection{Plano de Gestão de Configuração}
\subsection{Registro de Decisões Técnicas (ADRs)}
\subsection{Métricas de Acompanhamento}

\section{Capítulo 5 --- Processos Técnicos}
\subsection{Requisitos do Sistema}
\subsection{Arquitetura e Design}
\subsection{Implementação e Integração}
\subsection{Verificação, Validação e Transição}
\subsection{Operação, Manutenção e Descontinuação}

\end{document}
